% Page 015 — page imprimée 11
% Extrait documentaire : suite de « ROMAN (Jean), vers 1668 - 1715 »
% Mise en page en deux colonnes.


{\fontsize{9.65}{11.05}\selectfont
\begin{multicols}{2}

le fut enfermé à la Tour de Constance, en sortit en 1693, on ne sait pour quelle raison, et gagna en 1701 le Refuge à Londres (sa soeur Marguerite épousera un marchand réfugié à Dublin, en 1705).

\section*{Heures tragiques}

Roman regagna le Bougès : on sait par les aveux de l’une de ses anciennes fidèles qu’il avait des caches sûres dans quelques-unes des meilleures familles paysannes de la région du Bougès (hameaux de Saint-Julien d’Arpaon et La Salle-Prunet, ou chez les Velay (*), à Racoules). Au début de 1690, le prédicant se joignit à Vivent, avec quelques hommes armés, et apporta son aide à la préparation d’une insurrection générale des protestants, qu’étaient venus fomenter depuis la Suisse deux anciens pasteurs, Durand-Foncouverte et Dautun. L’entreprise avorta après la répression qui frappa l’importante assemblée de l’Espinas (région de Vialas), où s’étaient réunis les chefs du mouvement. Un an plus tard, le 10 mars 1691, Roman échappait de justesse à ses poursuivants, à Castagnols (Vialas), devait se jeter dans une rivière et marcher deux lieues en chemise, sous un gel cruel. En mars 1692, sa tête était mise à prix (300 livres), avec celle des autres prédicants; il se rendit à Genève, par Langogne et Le Puy, alors que l’abbé du Chaïla faisait courir le bruit calomnieux qu’il était indésirable dans la ville. Il en revint l’année suivante pour reprendre son apostolat errant. Il manqua être pris, après trahison, au hameau du Mijavols, que les soldats pillèrent (il l’avait été une première fois, après une assemblée, en mai 1686). En 1694, dans le même quartier, le prédicant fut blessé d’une balle à la jambe, et, transporté de caverne en caverne, manqua périr. Il était alors assisté dans son ministère par le jeune David Rey, de Massevaques (le hameau du futur chef camisard Castanet (*)). Il prêcha ensuite dans la région de Saint-Ambroix, dans celle de Meyrueis (où il avait quelques amis fidèles à Campis, et aussi un traître vainement attaché à sa perte), à Florac (il manqua être pris, les soldats fouillèrent le lit sous lequel il se cachait); il dut se résoudre à se faire accompagner de quelques hommes armés.

\section*{Pris et libéré}

En juin 1697, Roman repassait à Genève. Il y fut pris dans une querelle opposant, entre protestants, les modérés aux zélateurs, et fut traité par quelques-uns de «tentateur et perturbateur», cause par ses assemblées des malheurs qui frappaient les Réformés du Languedoc. Il se défendit vigoureusement, rétorquant à ses contradicteurs : «Présentement, il est question de savoir à qui Dieu redemandera ces âmes innocentes qui périssent faute de lumière»; Charles Bost rapproche cette controverse de celle qui opposait au même moment Brousson aux pasteurs de Hollande, et estime que les deux hommes «dans cet instant décisif (...) eurent l’âme pareille»; Roman, comme le nîmois, préféra venir retrouver son peuple en hautes Cévennes: il l’avait confié, dans les anciennes paroisses, à des «anciens» fidèles qui l’avaient entretenu dans la pratique et la foi. Il tint, le 29 décembre 1697, une très importante assemblée à Montcuq (Saint-Maurice de Ventalon), qui aurait regroupé 4 000 personnes, et fut suivie d’une forte répression, menée par l’abbé du Chaïla. L’insaisissable prédicant fut enfin pris, dans la nuit du 9 au 10 août 1699, près de Boucoiran, non loin de Nîmes. Trahi, gravement blessé, il fut enfermé sous bonne garde dans une auberge de Boucoiran; interrogé, il répondit en engageant des controverses. Mais à la foire de Lédignan, le 10 août, les protestants, informés de son arrestation, décidèrent de venir le libérer; une troupe d’hommes masqués vint assiéger l’auberge, et Roman fut délivré (il avait promis aux archers qui voulaient l’achever que les assaillants leur laisseraient la vie sauve, ce qui fut le cas). Roman quitta le village sur les épaules de ses frères, en entonnant le Psaume 51; de ceux qui le délivrèrent,

\end{multicols}
}


